\documentclass[11pt]{article}

\usepackage[utf8]{inputenc}
\usepackage[T1]{fontenc}
\usepackage[spanish,es-noquoting]{babel}
\usepackage{lmodern}
\usepackage[margin=2.5cm]{geometry}
\usepackage{enumitem}
\usepackage{titlesec}
\usepackage{parskip}

\setlist[itemize]{leftmargin=1.4em, itemsep=2pt, topsep=3pt}
\setlist[enumerate]{leftmargin=1.6em, itemsep=2pt, topsep=3pt}
\titleformat{\section}{\normalfont\large\bfseries}{\thesection.}{0.6em}{}
\titleformat{\subsection}{\normalfont\normalsize\bfseries}{\thesubsection}{0.6em}{}
\titlespacing{\section}{0pt}{12pt}{5pt}
\titlespacing{\subsection}{0pt}{9pt}{3pt}

\title{\vspace{-1.2cm}\textbf{Sistema de Gestión de Talleres Mecánicos}\\[4pt]
\large Propuesta de Idea — Grupo 10}
\author{Taller de Integración II \ $\cdot$ \ Taller de Integración IV}
\date{Universidad Católica de Temuco — Ingeniería Civil en Informática\\
Profesor: Ignacio Lincolao Venegas — Semestre II, 2026 — Metodología: Scrum}

\begin{document}
\maketitle
\tableofcontents
\bigskip

\section{Resumen ejecutivo}

Plataforma web y móvil para digitalizar y optimizar la gestión operativa de talleres mecánicos, desde la recepción del vehículo hasta su entrega al cliente.

Integra II desarrolla la aplicación web (frontend, backend y base de datos); Integra IV desarrolla la aplicación móvil. Ambos clientes comparten la misma API y base de datos, conformando un único sistema distribuido.

\section{El problema a resolver}

El sector de los talleres mecánicos opera desde hace años bajo dinámicas manuales, lo que genera una necesidad de modernización. Esta gestión tradicional produce:

\begin{itemize}
  \item Dificultad para reconstruir qué se le hace a cada vehículo, cuándo, con qué repuestos y por qué mecánico.
  \item Pérdida de información de diagnósticos, presupuestos y repuestos, al no existir un sistema centralizado.
  \item Mala comunicación con el cliente, quien al no saber el estado de su auto llama o escribe de forma constante.
  \item Interrupciones al mecánico, que muchas veces trabaja con las manos ocupadas y no puede contestar el teléfono.
  \item Escasa visibilidad operativa para el dueño del taller, que no puede ver de un vistazo el estado global del negocio.
\end{itemize}

\section{La solución}

Desarrollar una plataforma web y móvil que digitalice y optimice la gestión operativa del taller, desde que se recibe el vehículo hasta su entrega final.

Los trabajadores registran el ingreso de los autos, gestionan las horas agendadas y actualizan los avances de cada reparación. El cliente, por su parte, puede consultar el estado del servicio de su vehículo sin necesidad de contactar directamente al taller.

\section{Objetivos}

\subsection{Objetivo general}

Gestionar digitalmente el ciclo de reparación y mantención del vehículo, mediante una plataforma web y móvil que mejore el registro de la información, la comunicación con el cliente y la eficiencia operativa del taller.

\subsection{Objetivos específicos}

Comprometidos dentro del plazo del curso:

\begin{enumerate}
  \item Registrar vehículos asociados a un cliente, con marca, modelo, año, kilometraje y patente única.
  \item Gestionar el historial de mantenciones, reparaciones, repuestos utilizados y mecánico responsable.
  \item Implementar el flujo de estados de la orden, incluyendo los estados de reparación y los estados «Esperando aprobación de presupuesto» y «Esperando aprobación de documento».
  \item Permitir al cliente consultar el estado del servicio de su vehículo, con notificaciones y evidencia multimedia (fotografía o video).
  \item Registrar el distribuidor o proveedor asociado a cada repuesto, para que el cliente pueda verificar la procedencia y ver qué proveedor quedó asociado al repuesto que solicitó.
  \item Incorporar un módulo de control de inventario y stock que valide la disponibilidad de repuestos.
  \item Permitir configurar la cantidad de órdenes activas por mecánico y el número máximo de vehículos revisados por día.
  \item Diseñar una API REST centralizada compatible con la aplicación móvil de Integra IV.
  \item Implementar un panel de control con el estado global de los vehículos para el administrador.
  \item Ejecutar pruebas funcionales y de integración (unitarias, de API y de las actions) y desplegar la plataforma en un hosting con acceso público.
\end{enumerate}

\subsection{Fuera de plazo}

Quedan fuera del plazo comprometido: facturación electrónica, cotización automática a proveedores, pasarela de pagos, chat integrado, reserva de horas en línea y recordatorios avanzados de mantención.

\section{Actores del sistema}

\begin{itemize}
  \item \textbf{Cliente:} consultar el estado del servicio de su vehículo, revisar historial y ficha técnica por patente, aprobar presupuestos con evidencia y recibir notificaciones.
  \item \textbf{Mecánico:} actualizar el estado del vehículo, registrar diagnóstico y repuestos utilizados, adjuntar evidencia fotográfica o en video y consultar la ficha técnica del auto.
  \item \textbf{Administrador:} acceder al panel de control global, gestionar usuarios, consultar reportes, gestionar el inventario y configurar los límites operativos del taller.
\end{itemize}

\section{Flujo de estados del vehículo}

El flujo es secuencial y no puede saltarse etapas:

\begin{enumerate}
  \item Recibido
  \item Esperando Diagnóstico
  \item Esperando aprobación de presupuesto
  \item En Reparación
  \item Esperando Repuestos
  \item Esperando aprobación de documento
  \item Listo
  \item Entregado
\end{enumerate}

Cada transición queda registrada con fecha, hora y usuario responsable. En «Esperando aprobación de presupuesto» la orden no avanza hasta que el cliente aprueba; si lo rechaza, la orden queda en estado Rechazado. En «Esperando Repuestos» el sistema registra el repuesto solicitado y el distribuidor asociado, de modo que el cliente pueda conocer la procedencia de la pieza. En «Esperando aprobación de documento» se deja constancia del documento asociado a la entrega antes de marcar la orden como Listo.

\section{Propuesta de valor}

\subsection{Para el cliente}

\begin{itemize}
  \item Transparencia sobre el repuesto cambiado, su valor y el distribuidor del que proviene.
  \item Historial consultable por patente mediante ficha técnica.
  \item Posibilidad de consultar el estado del servicio, con notificaciones y presupuestos con evidencia fotográfica o en video.
\end{itemize}

\subsection{Para el taller y los mecánicos}

\begin{itemize}
  \item Menos interrupciones: al mantener informado al cliente mediante la app, el mecánico se concentra en su trabajo.
  \item Panel de control con el estado global de los vehículos para el administrador.
  \item Control de inventario y stock integrado a las órdenes de trabajo.
  \item Estandarización del trabajo mediante el flujo de estados definido para cada vehículo.
\end{itemize}

\section{Producto mínimo viable (MVP)}

\begin{enumerate}
  \item Registro e inicio de sesión con roles diferenciados.
  \item Registro de vehículos asociados a un cliente (marca, modelo, año, km, patente).
  \item Creación y gestión de órdenes de trabajo.
  \item Flujo completo de estados con historial de transiciones.
  \item Registro de repuestos, mecánico y distribuidor por intervención.
  \item Control de inventario y stock de repuestos.
  \item Notificaciones al cliente y soporte multimedia como evidencia.
  \item Historial de reparaciones consultable por patente.
  \item Panel de control del taller con estado global de vehículos.
  \item Parámetros configurables: órdenes activas por mecánico y máximo de vehículos revisados por día.
  \item API REST funcional para integración con la app móvil de Integra IV.
  \item Base de datos relacional implementada y desplegada en hosting con acceso público.
\end{enumerate}

\section{Arquitectura y stack tecnológico}

La arquitectura sigue un modelo cliente–servidor con API centralizada compartida entre la aplicación web y la móvil:

\begin{itemize}
  \item Web (Integra II): React.
  \item Móvil (Integra IV): Flutter o React Native.
  \item Backend: Python (API REST).
  \item Base de datos: PostgreSQL / MySQL.
  \item Autenticación: JWT.
  \item Versionamiento: Git / GitHub.
  \item Documentación API: OpenAPI / Swagger.
  \item Hosting: con acceso público, por definir con Integra IV.
\end{itemize}

\section{Modelo de datos preliminar}

Entidades principales y sus atributos:

\begin{itemize}
  \item \textbf{Usuario:} id, nombre, email, password\_hash, rol.
  \item \textbf{Cliente:} id, usuario\_id, teléfono.
  \item \textbf{Vehículo:} id, cliente\_id, patente, marca, modelo, año, km.
  \item \textbf{Mecánico:} id, usuario\_id, especialidad.
  \item \textbf{Orden de Trabajo:} id, vehículo\_id, mecánico\_id, estado, fechas.
  \item \textbf{Distribuidor:} id, nombre, RUT, contacto, es\_oficial.
  \item \textbf{Repuesto:} id, nombre, precio, stock.
  \item \textbf{Detalle Reparación:} id, orden\_id, repuesto\_id, distribuidor\_id, cantidad.
  \item \textbf{Historial de Estado:} id, orden\_id, estado\_anterior, estado\_nuevo, fecha.
  \item \textbf{Presupuesto:} id, orden\_id, monto, estado, url\_evidencia.
\end{itemize}

\section{Metodología — Scrum}

El proyecto se desarrolla con la metodología ágil Scrum. Roles:

\begin{itemize}
  \item \textbf{Scrum Master:} facilitar eventos y eliminar impedimentos.
  \item \textbf{Product Owner:} gestionar el backlog y priorizar historias de usuario.
  \item \textbf{Dev Team:} desarrollar los módulos del sistema.
\end{itemize}

Eventos: Sprint Planning, Daily Stand-up (10–15 min), Sprint Review y Retrospectiva.

\subsection{Evaluaciones del curso}

\begin{itemize}
  \item Sistematización del proyecto: 20\%.
  \item Presentación 1.\textsuperscript{er} Sprint: 20\%.
  \item Presentación 2.\textsuperscript{do} Sprint: 20\%.
  \item Sprint Final: 40\%.
\end{itemize}

\section{Riesgos identificados}

\begin{itemize}
  \item \textbf{Alcance excesivo} (impacto alto): MVP bien definido y backlog priorizado.
  \item \textbf{Desalineación con Integra IV} (alto): contrato de API acordado antes de desarrollar.
  \item \textbf{Modelo de datos incompleto} (alto): validar el ERD en las primeras semanas.
  \item \textbf{Integración web–móvil} (medio): lógica centralizada en el backend.
  \item \textbf{Autenticación y roles} (medio): definir roles y permisos desde el diseño.
\end{itemize}

\section{Reglas de negocio}

\subsection{Vehículos y clientes}
\begin{itemize}
  \item Un vehículo debe estar asociado a un cliente registrado. La patente es única e irrepetible.
  \item Un cliente puede tener múltiples vehículos, pero solo consulta los suyos.
\end{itemize}

\subsection{Órdenes de trabajo}
\begin{itemize}
  \item Un vehículo solo puede tener una orden activa a la vez.
  \item No se puede crear una orden sin mecánico asignado.
  \item La fecha de ingreso se registra automáticamente y no se puede modificar.
  \item Solo el administrador puede eliminar una orden.
  \item La cantidad máxima de órdenes activas por mecánico es configurable por el administrador.
\end{itemize}

\subsection{Flujo de estados}
\begin{itemize}
  \item El flujo es secuencial y no puede saltarse etapas.
  \item Cada cambio de estado queda registrado con fecha, hora y usuario.
  \item Solo mecánico o administrador puede cambiar el estado.
  \item El estado Entregado es irreversible.
  \item «En Reparación» requiere que el presupuesto haya sido aprobado por el cliente.
  \item «Esperando Repuestos» requiere al menos un repuesto y distribuidor asociado.
\end{itemize}

\subsection{Repuestos, distribuidores e inventario}
\begin{itemize}
  \item Todo repuesto debe tener un distribuidor asociado, marcado como oficial o genérico.
  \item El precio del repuesto no puede modificarse una vez aprobado el presupuesto.
  \item No se puede registrar un repuesto como utilizado si no hay stock disponible.
  \item El stock se descuenta al registrar el consumo de un repuesto en una orden.
\end{itemize}

\subsection{Presupuestos}
\begin{itemize}
  \item No se puede avanzar a «En Reparación» sin presupuesto aprobado por el cliente.
  \item Todo presupuesto debe incluir al menos una foto o video como evidencia.
  \item Si el cliente rechaza el presupuesto, la orden queda bloqueada en estado Rechazado.
\end{itemize}

\subsection{Configuración operativa}
\begin{itemize}
  \item El administrador puede configurar la cantidad de órdenes activas permitidas por mecánico.
  \item El administrador puede establecer el número máximo de vehículos revisados por día.
\end{itemize}

\subsection{Historial}
\begin{itemize}
  \item Se mantiene un historial con cada cambio, reparación, cliente o taller involucrado.
  \item El historial es de solo lectura: nadie puede modificarlo una vez registrado.
  \item El historial de un vehículo se conserva aunque el cliente elimine su cuenta.
\end{itemize}

\section{Requisitos funcionales}

\subsection{Gestión de usuarios y acceso}
\begin{itemize}
  \item RF01: registrar usuarios con nombre, email y contraseña.
  \item RF02: autenticar usuarios mediante email y contraseña.
  \item RF03: diferenciar tres roles: Cliente, Mecánico y Administrador.
  \item RF04: el administrador puede crear, editar y desactivar cuentas de mecánicos.
  \item RF05: cada usuario ve únicamente la información correspondiente a su rol.
\end{itemize}

\subsection{Gestión de vehículos}
\begin{itemize}
  \item RF06: registrar vehículos asociados a un cliente, con patente, marca, modelo, año y kilometraje.
  \item RF07: validar que no existan dos vehículos con la misma patente.
  \item RF08: el cliente puede consultar la ficha técnica de su vehículo por patente.
  \item RF09: un cliente puede registrar múltiples vehículos.
\end{itemize}

\subsection{Órdenes de trabajo}
\begin{itemize}
  \item RF10: crear órdenes asociadas a un vehículo y un mecánico.
  \item RF11: registrar automáticamente la fecha de ingreso al crear la orden.
  \item RF12: un vehículo solo puede tener una orden activa a la vez.
  \item RF13: solo el administrador puede eliminar una orden.
\end{itemize}

\subsection{Flujo de estados}
\begin{itemize}
  \item RF14: gestionar el flujo Recibido → Esperando Diagnóstico → Esperando aprobación de presupuesto → En Reparación → Esperando Repuestos → Esperando aprobación de documento → Listo → Entregado.
  \item RF15: registrar cada cambio de estado con fecha, hora y usuario responsable.
  \item RF16: solo mecánicos y administradores pueden cambiar el estado.
  \item RF17: el estado Entregado es irreversible.
\end{itemize}

\subsection{Presupuestos y evidencia}
\begin{itemize}
  \item RF18: crear presupuestos con monto y evidencia fotográfica o en video.
  \item RF19: el cliente puede aprobar o rechazar el presupuesto desde su vista.
  \item RF20: no se puede avanzar a «En Reparación» sin presupuesto aprobado.
  \item RF21: si el cliente rechaza el presupuesto, la orden queda en estado Rechazado.
\end{itemize}

\subsection{Repuestos, distribuidores e inventario}
\begin{itemize}
  \item RF22: registrar repuestos con nombre, precio y stock.
  \item RF23: todo repuesto utilizado en una orden debe tener un distribuidor asociado.
  \item RF24: el distribuidor se marca como oficial o genérico.
  \item RF25: validar el stock disponible antes de registrar un repuesto como utilizado.
  \item RF26: descontar el stock al registrar el consumo de un repuesto.
  \item RF27: el cliente puede ver qué distribuidor o proveedor quedó asociado al repuesto de su orden.
\end{itemize}

\subsection{Notificaciones y multimedia}
\begin{itemize}
  \item RF28: notificar al cliente en cada cambio de estado de su orden.
  \item RF29: permitir adjuntar archivos multimedia (fotografías o video) como evidencia en presupuestos y reparaciones.
\end{itemize}

\subsection{Configuración operativa}
\begin{itemize}
  \item RF30: permitir configurar la cantidad de órdenes activas por mecánico.
  \item RF31: permitir establecer el número máximo de vehículos revisados por día.
\end{itemize}

\subsection{Panel de control}
\begin{itemize}
  \item RF32: el administrador dispone de un panel con el estado global de todos los vehículos en el taller.
  \item RF33: el mecánico ve únicamente las órdenes que le fueron asignadas.
\end{itemize}

\subsection{Historial}
\begin{itemize}
  \item RF34: mantener un historial con cada cambio, reparación, cliente o taller.
  \item RF35: el historial es de solo lectura, sin posibilidad de edición.
  \item RF36: el historial permite reconstruir la trayectoria de un vehículo de forma prácticamente inmediata, en cuestión de segundos.
  \item RF37: el historial se conserva aunque el cliente elimine su cuenta.
\end{itemize}

\subsection{Integración con la app móvil}
\begin{itemize}
  \item RF38: exponer una API REST documentada para la integración con la app móvil de Integra IV.
  \item RF39: los cambios realizados desde la app móvil se reflejan en la web.
\end{itemize}

\section{Pruebas y despliegue}

\subsection{Pruebas}
\begin{itemize}
  \item Pruebas funcionales de cada módulo del sistema.
  \item Pruebas de integración entre el backend y la aplicación móvil.
  \item Pruebas unitarias sobre la lógica de negocio.
  \item Pruebas sobre los endpoints de la API.
  \item Pruebas de las actions (flujos de trabajo automatizados) del proyecto.
\end{itemize}

\subsection{Despliegue}
\begin{itemize}
  \item Despliegue de la aplicación y de la base de datos en un hosting con acceso público.
\end{itemize}

\section{Requisitos no funcionales}

\subsection{Seguridad}
\begin{itemize}
  \item RNF01: todas las comunicaciones se realizan bajo HTTPS.
  \item RNF02: protección de datos personales de clientes y trabajadores; las contraseñas se almacenan con hash seguro, nunca en texto plano.
  \item RNF03: la autenticación se implementa mediante JWT con expiración de sesión.
  \item RNF04: el backend valida la autorización en cada endpoint.
  \item RNF05: un usuario no puede acceder a recursos de otro usuario.
\end{itemize}

\subsection{Rendimiento}
\begin{itemize}
  \item RNF06: el sistema responde a las solicitudes en un tiempo acotado bajo condiciones normales de uso.
  \item RNF07: la API soporta múltiples usuarios concurrentes sin degradación relevante del servicio.
\end{itemize}

\subsection{Usabilidad y compatibilidad}
\begin{itemize}
  \item RNF08: la interfaz es intuitiva y responsiva, y funciona en dispositivos móviles y de escritorio.
  \item RNF09: los mensajes de error son claros y orientan al usuario.
  \item RNF10: la aplicación web funciona en las versiones actuales de los navegadores más usados.
\end{itemize}

\subsection{Mantenibilidad y escalabilidad}
\begin{itemize}
  \item RNF11: el código se versiona en GitHub con commits descriptivos.
  \item RNF12: la API se documenta con OpenAPI/Swagger.
  \item RNF13: la arquitectura permite agregar módulos sin reescribir el sistema base.
\end{itemize}

\section{Conclusión}

El Sistema de Gestión de Talleres Mecánicos aborda un problema concreto y frecuente en talleres de pequeño y mediano tamaño, cubriendo tanto las necesidades operativas del taller como la experiencia del cliente.

La propuesta se apoya en un MVP acotado, una lógica de backend clara y una arquitectura escalable, integrada con la aplicación móvil de Integra IV y verificada mediante pruebas funcionales y de integración.

\end{document}
